% Fourier-tail and multilinear-jet estimates used in the regular-source chart.
\section{Fourier tails and local jet estimates}
\label{sec:ex-lattice-trace-row-jet-estimates}

This section proves two quantitative ingredients of the small-support
regular-source argument: a weighted convolution bound with an exponentially
small tail, and the corresponding multilinear estimate for a local
coefficient probe. The remaining hypotheses are named mathematically: a
uniform Fourier strip bound, a coefficient-growth receipt, and the local
regularity and vanishing-jet receipts. In the construction these come from
the selected finite-range kernel, the budget's order bounds, and the
regularity of the masked source.

For each natural number \(n\), write \(\mathbb Z^4\) for the covering lattice, with
\(|x|_1=\sum_{\nu=1}^4|x_\nu|\), and put
\[
  W_n(x)=(1+|x|_1)^n,\qquad
  C_n(\sigma)=\left(1+\frac{n+1}{\sigma}\right)^n
  \quad(\sigma>0).
 \]
The elementary inequality
\[
 (1+r)^n\le C_n(\sigma)e^{\sigma r}\qquad(r\ge0)
 \tag{1}\label{eq:ex-lattice-widthcost}
 \]
follows by setting \(s=\sigma/(n+1)\). Indeed
\[
 1+r\le\left(1+\frac{n+1}{\sigma}\right)(1+sr)
       \le\left(1+\frac{n+1}{\sigma}\right)e^{sr},
 \]
and taking the \(n\)-th power gives (1), since \(ns\le\sigma\).
Also \(W_n(x+z)\le W_n(x)W_n(z)\), by the triangle inequality and
\(1+a+b\le(1+a)(1+b)\).

\begin{lemma}[Polynomially weighted convolution]
\label{lem:ex-lattice-row-convolution}
Suppose a scalar kernel \(q:\mathbb Z^4\to\mathbb R\) and constants
\(\sigma>0\), \(C\ge0\) satisfy
\[
 \sum_{z\in\mathbb Z^4}e^{\sigma|z|_1}|q(z)|\le C
 \tag{2}\label{eq:ex-lattice-row-strip}
 \]
and a Banach-valued array \(b\) satisfies
\(\|b(x)\|\le M W_n(x)\), where \(M\ge0\). Then the convolution
\[
 (q*b)(x)=\sum_z q(z)b(x+z)
 \]
is absolutely convergent and
\[
 \|(q*b)(x)\|\le
 M C_n(\sigma) C W_n(x).
 \tag{3}\label{eq:ex-lattice-row-convolution}
 \]
For a \(4\)-by-\(4\) matrix kernel whose sixteen entries each satisfy
(2) with the same \(C\), the corresponding vector-row operator obeys
\[
 \|(\mathsf A b)(x)\|\le
 4M C_n(\sigma) C W_n(x),
 \tag{4}\label{eq:ex-lattice-row-matrix-cost}
 \]
when the output norm is the maximum of its four component norms.
\end{lemma}
\begin{proof}
For each \(z\),
\[
 |q(z)|\,\|b(x+z)\|
 \le M W_n(x)|q(z)|W_n(z)
 \le M W_n(x)C_n(\sigma)e^{\sigma|z|_1}|q(z)|,
 \]
where (1) was applied to \(|z|_1\). Summing proves absolute convergence and
(3). For a matrix row, apply (3) to each of its four entries and sum the
four resulting component bounds; this gives (4).
\end{proof}

\begin{lemma}[Agreement on a ball gives an exponential tail]
\label{lem:ex-lattice-row-tail}
Assume (2). Let \(A_1,A_2:\mathbb Z^4\to E\), where \(E\) is a Banach
space, satisfy
\[
 \|A_1(y)\|\le d_1W_n(y),\qquad
 \|A_2(y)\|\le d_2W_n(y),
 \]
with \(d_1,d_2\ge0\), and suppose \(A_1(x+z)=A_2(x+z)\) whenever
\(|z|_1<b\). Then
\[
 \left\|\sum_z q(z)\bigl(A_1(x+z)-A_2(x+z)\bigr)\right\|
 \le e^{-\sigma b/2}(d_1+d_2)C_n(\sigma/2)C W_n(x).
 \tag{5}\label{eq:ex-lattice-row-tail}
 \]
\end{lemma}
\begin{proof}
The difference vanishes for \(|z|_1<b\), so the triangle inequality and
submultiplicativity of \(W_n\) give an upper bound
\[
 (d_1+d_2)W_n(x)\sum_{|z|_1\ge b}|q(z)|W_n(z).
 \]
Using (1) with \(\sigma/2\), each summand is at most
\[
 C_n(\sigma/2)e^{\sigma|z|_1/2}|q(z)|
 \le C_n(\sigma/2)e^{-\sigma b/2}
       e^{\sigma|z|_1}|q(z)|.
 \]
Summation and (2) prove (5).
\end{proof}

\begin{lemma}[Multilinear jet comparison]
\label{lem:ex-lattice-jet-comparison}
Here is the passage from that tail estimate to a local coefficient
measurement. Let \(P\) be a linear screen onto the allowed edge
directions; it is applied only to the direction arrays below. Let \(F\) be
a \(C^N\) real-valued local observable. Here \(\mathrm{inc}_1(U)\) is the
configuration obtained by exponentiating the tangent field \(U\) from the
identity. Write
\(\mathcal J_a=D^a(F\circ\mathrm{inc}_1)(0)\) for its symmetric
\(a\)-linear derivative at the identity configuration. Assume the following
three estimates for \(a\le N\):
\begin{enumerate}
\item for all admissible screened directions \(v_1,\ldots,v_a\), and all
nonnegative \(\lambda_i\) with \(\|v_i\|\le\lambda_i\),
\[
 |\mathcal J_a(v_1,\ldots,v_a)|
 \le \Phi w\prod_{i=1}^a\lambda_i;
 \tag{6}\label{eq:ex-lattice-jet-price}
\]
\item the already-screened reference directions \(A_i=P(\theta y_i)\)
satisfy \(\|A_i\|\le K\), with \(K\ge1\), and the screened reference jet
vanishes:
\[
 \mathcal J_a(P(\theta y_1),\ldots,P(\theta y_a))=0;
 \tag{7}\label{eq:ex-lattice-jet-vanishing}
\]
\item directions \(x_i\) approximate them with
\(\|P(\theta x_i)-P(\theta y_i)\|\le\varepsilon\), where
\(0\le\varepsilon\le K\).
\end{enumerate}
Then
\[
 |\mathcal J_a(P(\theta x_1),\ldots,P(\theta x_a))|
 \le 2^a\varepsilon K^a\Phi w.
 \tag{8}\label{eq:ex-lattice-jet-comparison}
\]
\end{lemma}
\begin{proof}
Set \(U_i=P(\theta x_i)\), \(A_i=P(\theta y_i)\), and
\(E_i=U_i-A_i\). By multilinearity,
\[
 \mathcal J_a(U_1,\ldots,U_a)-\mathcal J_a(A_1,\ldots,A_a)
 =\sum_{\varnothing\ne S\subseteq\{1,\ldots,a\}}
 \mathcal J_a(v_1^S,\ldots,v_a^S),
 \]
where \(v_i^S=E_i\) if \(i\in S\), and \(v_i^S=A_i\) otherwise.
The empty-set term is zero by (7). In every remaining term at least one
factor has norm at most \(\varepsilon\); all other factors have norm at
most \(K\), by the stated bounds on the screened reference directions.
Applying (6) to each term bounds it by
\(\varepsilon K^{a-1}\Phi w\le\varepsilon K^a\Phi w\).
There are \(2^a-1\le2^a\) such terms, proving (8).
\end{proof}

\begin{proposition}[Local probe bound from a strip and a jet gap]
\label{prop:ex-lattice-strip-to-probe}
Fix natural numbers \(d,N\) with \(d\le N\), a support radius \(b\), a
chart radius \(m\), and \(\sigma>0\), and let \(L\ge1\) be the integer
block side. For each
\(k\le d\), suppose the four-by-four row kernel
\((q_{\nu j})_{\nu,j=1}^4\) has entries satisfying (2) with constant
\(M\). Suppose the two arrays \(A^{(0)}_j,A^{(1)}_j\) satisfy
\[
 \|A^{(h)}_j(y)\|\le K W_{dk}(y)\quad(h=0,1),
 \qquad
 A^{(0)}_j(u+z)=A^{(1)}_j(u+z)\ \text{if }|z|_1<b,
 \tag{9a}\label{eq:ex-lattice-local-array-input}
\]
where \(K\ge1\), \(M\ge0\), and each kernel base \(u\) has
\(|u|_1\le4m\). Define the order-\(k\) reference and actual directions \(y_k\) and
\(x_k\) by
applying this matrix row to \(A^{(0)}\) and \(A^{(1)}\), respectively.
The screen \(P\) is coordinate restriction to allowed edges. Let
\(\theta>0\) be the threshold and \(\omega(e)\ge1\) the spatial weight of
the budget; the norm in (6) is the normalized edge norm
\[
 \|v\|=\sup_{e}\frac{\|v(e)\|}{\theta\,\omega(e)},
 \qquad\text{so that}\qquad
 \|P(\theta v)\|\le\sup_{e\ \mathrm{allowed}}\|v(e)\|
\]
for every direction array \(v\). Thus an entrywise bound for an unscaled
direction on the allowed edges is a bound for the normalized norm of its
threshold-scaled screen, and no additional screen norm is inserted
between the row estimate and the priced jet. Then the entrywise
difference of \(x_k\) and \(y_k\) is at most
\[
 \delta_k=
  8K e^{-\sigma b/2}(1+4m)^{dk}C_{dk}(\sigma/2)M
 \tag{9}\label{eq:ex-lattice-local-delta}
\]
and hence so is the normalized norm of \(P(\theta x_k)-P(\theta y_k)\).
The reference and actual directions satisfy, entrywise,
\[
 \|y_k\|,\ \|x_k\|
 \le 4K C_{dk}(\sigma)M(1+4L(m+2))^{dk}\le C_1,
\]
where
\[
 C_1=1+\sum_{k=0}^{d}
 4K C_{dk}(\sigma)M\bigl(1+4L(m+2)\bigr)^{dk},
 \tag{10}\label{eq:ex-lattice-local-C1}
\]
Assume the local jet satisfies the price bound (6), that \(F\) depends
only on the allowed-edge components of the configuration, so that
\(\mathcal J_a(v_1,\ldots,v_a)=\mathcal J_a(Pv_1,\ldots,Pv_a)\), and also
that
\(\varepsilon_d\le C_1\). Assume the weighted vanishing receipt for the
screened, threshold-scaled reference directions:
\[
 \mathcal J_a\bigl(P(\theta y_{q_1}),\ldots,P(\theta y_{q_a})\bigr)=0
 \quad\text{whenever }q_i\ge1,\quad \sum_{i=1}^a q_i\le d.
 \tag{10a}\label{eq:ex-lattice-weighted-jet-vanishing}
\]
In the application, (10a) holds because each \(P(\theta y_k)\) lies in the
space \(\mathcal V_k\) of polynomial directions of degree at most \(k-1\)
on the allowed edges, and the observable has vanishing jets on lists of
such directions with positive degrees of sum at most \(d\). No such
membership is assumed for the actual directions \(x_k\); it is neither
needed below nor available in the application.
Assume also the order-zero receipt \(F(\mathrm{id})=0\).
Let \(\gamma:\mathbb R\to\mathcal X\) be the actual \(C^d\) probe path in
the tangent configuration space, with \(\gamma(0)=0\) and
\(\gamma^{(k)}(0)=\theta x_k\) for \(1\le k\le d\). The observable is
evaluated on
\(\operatorname{inc}_1(\gamma(s))\), so its base point is the identity
configuration. Then the probe of total derivative order \(\ell\le d\)
obeys
\[
 \left|\left.\frac{d^\ell}{ds^\ell}
 F(\operatorname{inc}_1(\gamma(s)))\right|_{s=0}\right|
 \le |\operatorname{Part}^{\mathrm{ord}}(\ell)|\,2^\ell
       \varepsilon_d C_1^\ell\Phi w,\qquad
 \varepsilon_d=\sum_{k=0}^{d}\delta_k.
 \tag{11}\label{eq:ex-lattice-local-probe}
\]
\end{proposition}
\begin{proof}
For each of the four entries in an output row, (5) with
\(d_1=d_2=K\) bounds the tail by
\[
 2K e^{-\sigma b/2}C_{dk}(\sigma/2)M W_{dk}(u).
\]
There are four input components, so the row difference is at most
\(8K e^{-\sigma b/2}C_{dk}(\sigma/2)M W_{dk}(u)\).
Since \(|u|_1\le4m\), this is (9). Summing over \(k\le d\) gives a
common error bound \(\varepsilon_d\). For either direction, (4) and the
growth bound on its input array give
\(4K C_{dk}(\sigma)M W_{dk}(u)\). In the chart the lifted base index has
size at most \(4L(m+2)\), giving the order-\(k\) bound above; the extra
\(1\) ensures \(C_1\ge1\).

Here \(\operatorname{Part}^{\mathrm{ord}}(\ell)\) is the finite index set of
ordered finite partitions of the \(\ell\) differentiation slots into
nonempty blocks, in the canonical ordering used by the chain rule; it is
not a sum over arbitrary permutations of each set partition. For (11), use
the repeated chain rule in that ordered-finpartition form. When \(\ell=0\),
the derivative is \(F(\mathrm{id})=0\) by the separate order-zero receipt.
For each nonempty-order partition \(\pi\) of the \(\ell\) probe derivatives, the
chain-rule term is the \(|\pi|\)-linear jet \(\mathcal J_{|\pi|}\) applied
to the path derivative \(\theta x_{|B|}\) for each block \(B\in\pi\),
and by the locality assumption it is unchanged when each
\(\theta x_{|B|}\) is replaced by \(P(\theta x_{|B|})\).
The all-reference term, with \(P(\theta x_{|B|})\) replaced by
\(P(\theta y_{|B|})\), vanishes:
the block degrees are positive and sum to \(\ell\le d\), so the weighted
vanishing hypothesis applies. In the normalized norm, each screened
reference direction \(P(\theta y_{|B|})\) is bounded by
\(C_1\), and each \(P(\theta x_{|B|})\) differs from it by at most
\(\delta_{|B|}\le\varepsilon_d\). Since \(\varepsilon_d\le C_1\), the multilinear
comparison lemma (8), with \(K=C_1\), bounds this chain-rule term by
\(2^{|\pi|}\varepsilon_d C_1^{|\pi|}\Phi w\), which is at most
\(2^\ell\varepsilon_d C_1^\ell\Phi w\). Summing over the finite ordered
finpartition index set proves (11).
\end{proof}

\paragraph{The arrays in the selected construction.}
In the selected construction the probe is the composition of the local
observable with a routed test path: the logarithmic-link test curve of the
coefficient probe is cropped to the chart of radius \(m\) and passed
through the first-order row operator of the step. This routed path is
\(\theta\) times a real-analytic path \(R(s)\) of direction arrays with
\(R(0)=0\), and the actual direction is \(x_k=R^{(k)}(0)\). The two input
arrays, before the row operator is applied, are the following. The
reference array \(A^{(0)}=W_k\) is the \(k\)-th derivative at \(s=0\)
of the test script, an array on the covering lattice; the coefficient-growth
receipt is \(\|W_k(y)\|\le KW_{dk}(y)\). The actual array is its patched
periodic lift
\[
 A^{(1)}(z)=V_k(z)=
 \begin{cases}
  W_k(\operatorname{lift}(\operatorname{wrap}z)),
     &\operatorname{sep}(0,\operatorname{wrap}z)\le m,\\
  0,&\text{otherwise},
 \end{cases}
\]
where \(\operatorname{wrap}\) is the projection of the covering lattice
onto the periodic coarse lattice and \(\operatorname{lift}\) its centred
section; it satisfies the same growth bound. The placement hypotheses
stated below imply \(V_k(u+z)=W_k(u+z)\) for \(|z|_1<b\) at every kernel
base \(u\) that occurs, with \(|u|_1\le4m\). Applying the same row operator
to \(W_k\) and \(V_k\) gives the reference direction \(y_k\) and the actual
direction \(x_k\). The input arrays agree on the ball of radius \(b\);
the two convolved
directions need not agree anywhere. The reference direction is, on the
allowed edges, a polynomial in the spatial coordinates of degree at most
\(k-1\) when \(k\ge1\); this polynomial property is established for the
reference direction only, and is the membership
\(P(\theta y_k)\in\mathcal V_k\) used for (10a). The strip bound supplies
(2), while the
coefficient-growth receipt supplies \(K W_{dk}\). Thus the difference is
not estimated term by term without cancellation: all terms inside the
agreement ball cancel, and only the exponentially weighted exterior tail
is paid.

The hypotheses under which the source applies
Proposition~\ref{prop:ex-lattice-strip-to-probe} in this way are the
following; none of them is proved in this section.
\begin{enumerate}
\item The row operator applied to both arrays is literally the first-order
row of the actual step. It satisfies the strip bound (2) with constants
\(\sigma>0\), \(M\ge0\); it is identical to the row of the selected
kernel, and the patch block side of the step equals that
kernel's block side \(L\). A strip bound for some other kernel does not
suffice.
\item Placement: with \(c\) the side of the coarse periodic lattice,
\(20(m+b+1)<c\); the support \(X\) of the observable and the chart region
\(Y\) satisfy that the radius-\(b\) halo of the parent of every site of
\(X\) is contained in \(Y\), and every site of \(Y\) has separation at most
\(m\) from the origin.
\item The observable is regular on \(X\) for the budget, with all ordered
derivatives of order at most the budget's order \(n_B\) along normalized
directions bounded by \(w\) on healthy configurations; this yields (6)
with \(\Phi\) the fixed jet price of order \(N\), where
\(d\le N\le n_B\).
\item The observable has the polynomial-jet vanishing property of order
\(d\) on \(X\): its jets at the identity vanish on every list of directions
whose screens lie in \(\mathcal V_{q_1},\ldots,\mathcal V_{q_a}\) with
\(q_i\ge1\) and \(\sum_iq_i\le d\). This includes \(F(\mathrm{id})=0\) as
the case of the empty list.
\item The growth receipt \(\|W_k(y)\|\le KW_{dk}(y)\) holds for every
\(k\le d\) with one \(K\ge1\), the two integer parameters \((p,u)\) of
the selected kernel satisfy \(p\ge d\) and \(u\ge4\), its real parameter is
positive, and \(\varepsilon_d\le C_1\).
\end{enumerate}
Under these hypotheses the routed-path identity gives the path \(\gamma=
\theta R\), the tail lemma gives the entrywise error \(\delta_k\) on the
allowed edges, the row bound (4) gives \(\|y_k\|\le C_1\) on those edges,
and the polynomial property of \(y_k\) gives (10a); the proposition then
yields (11) with \(L\) the common block side.

The argument above proves the Fourier-tail and multilinear comparison
once the three receipts (strip, growth, and vanishing priced jet) are
available. It does not prove those upstream receipts, nor does it by
itself prove the source-placement or healthy-chart hypotheses in the
small-support branch. Those are separate construction inputs.
